\section{Adding methods and preserving comparable versions}
\label{app:rules}

The complete rules are in \texttt{benchmark/RULES.md}. The public runner
supports a 38-parameter wing, three tasks, a 4{,}096-surrogate cap and a
32-query CFD horizon. It records proposals, feedback and private nominations;
\texttt{verify.py} checks the trace and \texttt{score.py} summarizes outcomes.
The historical studies' actual access conditions and exceptions are reported
in Section~\ref{sec:analysis} and Appendix~\ref{app:resources}; they must not be
inferred from the public interface's maximum allowances.

\subsection{Integration contract}

A new method subclasses \texttt{Optimizer} in
\texttt{benchmark/optimizer\_api.py}. Its three decision methods are
\texttt{propose(ctx)}, returning the next 38-coordinate wing;
\texttt{observe(result)}, receiving the score or failure codes; and
\texttt{recommend()}, returning its current recommendation. Optional
\texttt{start(ctx)} initializes state. The context supplies the task, design
box, history, remaining budgets and the surrogate query function. A method
may maintain local state and call a language model inside this interface;
the evaluator and private grading contract need not change.

The worked optimizer and recorded-result evaluator in
\texttt{benchmark/README.md} exercise integration without CFD. The recorded
backend can answer only already graded wings and cannot validate a new
optimizer's aerodynamic performance. Fresh evaluations use the CFD backend
and the documented container. The original distributed agent campaign
controller is not included; reproducing that campaign's orchestration
requires additional engineering. The released interface, prompts and tools
support new implementations, not one-command reproduction of that controller.

\subsection{Versioning and maintenance scope}

The preservation policy is to retain the original records and content-pinned
protocol, distinguish documentation corrections from evaluator changes, and
report results under different mesh, solver, budget or outcome-policy versions
separately. A rebuilt runtime receives its own fingerprint and protocol
revision; it is not silently substituted for the original environment.
The container recipe, source/weight hashes, recorded example and analysis
scripts provide checks against dependency drift. Issues and integration
instructions are maintained with the repository, but the release offers no
hosted grading service or guaranteed maintenance duration. These mechanisms
reduce dependence on a live service; they do not guarantee that dependencies
will remain available or that future maintainers will have CFD resources.

For future submissions, the runner's closed track permits the provided
source data and warm-start records; an open track also permits public
benchmark records. The historical agents did not receive the 62-record
baseline context. Agent seed labels identify episodes and do not create
paired random draws across methods. Report actual resource use and protocol
version alongside recommendation validity and mean $L/D$ among valid
recommendations. A failed computation is a failure under that protocol,
not a certificate of physical infeasibility.
